
\section{Model Architecture Details}\label{app:architecture}

All our experiments use a single VLM architecture template, parameterised across the four scales of the scaling ladder (\cref{tab:scales}). The template follows the InternVL-3~\citep{zhu2025internvl3} family: a vision encoder $\rightarrow$ a randomly-initialised MLP projector $\rightarrow$ an autoregressive language-model backbone, with all three components trained jointly from the start (single-stage pretraining, no frozen components). Across our four scales, only the language-model backbone changes; the vision encoder and the projector recipe are held fixed. We document each component in turn.

\subsection{Vision Encoder}\label{app:arch-vision}

We use \textbf{InternViT-300M-448px-V2.5}~\citep{chen2024expanding} for all experiments. It is a Vision Transformer~\citep{dosovitskiy2020image} with the modifications introduced in the InternVL series~\citep{chen2024internvl,chen2024expanding,zhu2025internvl3}, kept identical across our four scales. \cref{tab:arch-vision} reports its key structural choices.

\begin{table}[!h]
\centering
\caption{\textbf{Vision encoder architecture (InternViT-300M-448px-V2.5).} The same vision encoder is used across all four scales. 
The encoder is fully unfrozen and updated jointly with the projector and the LM backbone.}
\label{tab:arch-vision}
\small
\begin{tabular}{ll}
\toprule
\textbf{Component} & \textbf{Value} \\
\midrule
\rowcolor{textcolor}\multicolumn{2}{l}{\textit{Input}} \\
Image resolution & $448 \times 448$ (one tile; dynamic high-res tiling, see \cref{sec:methodology}) \\
Patch size & $14 \times 14$ \\
Tokens per tile (pre-pixel-shuffle) & $32 \times 32 = 1024$ \\
Tokens per tile (post-pixel-shuffle, fed to LM) & $16 \times 16 = 256$ \\
Channels & 3 (RGB) \\
\midrule
\rowcolor{textcolor}\multicolumn{2}{l}{\textit{Transformer trunk}} \\
Depth (layers) & 24 \\
Hidden size & 1024 \\
Attention heads & 16 \\
Head dim & 64 \\
FFN intermediate size & 4096 \\
FFN activation & GELU \\
FFN style & 2-layer MLP (Linear $\rightarrow$ GELU $\rightarrow$ Linear) \\
Attention style & Standard multi-head; QKV bias \cmark, O-proj bias \xmark \\
QK normalisation & \xmark \\
Normalisation & Pre-LayerNorm, $\varepsilon = 10^{-6}$ \\
Positional embeddings & Learned absolute (interpolated to 448px) \\
Flash attention & \cmark (FA-2~\citep{dao2023flashattention}) \\
\midrule
\rowcolor{textcolor}\multicolumn{2}{l}{\textit{Total}} \\
Parameters & $\sim$304M \\
\bottomrule
\end{tabular}
\end{table}

The tile-based tokenisation produces $32 \times 32 = 1024$ patches per $448 \times 448$ tile. After the projector's pixel-shuffle reduction (\cref{app:arch-projector}), this becomes $16 \times 16 = 256$ visual tokens per tile ($4{\times}$ reduction) that are fed to the language model. 
Multiple tiles and a thumbnail image are concatenated into the LM input, following the dynamic high-resolution scheme of InternVL-2.5~\citep{chen2024expanding}.

\subsection{Projector}\label{app:arch-projector}

The vision encoder and the language model are bridged by a small randomly-initialised MLP-style projector (often called the ``connector'' in VLM literature~\citep{mckinzie2024mm1,liu2023visual}). It is the only module that is randomly initialised at the start of training, everything else is loaded from pretrained checkpoints. \cref{tab:arch-projector} gives the exact structure.
\begin{table}[!h]
\centering
\caption{\textbf{Projector architecture.} The projector is a fixed-depth, fixed-activation 2-layer MLP whose width is the only quantity that varies across scales (it tracks $D_{\mathrm{LM}}$, the language-model hidden size).}
\label{tab:arch-projector}
\small
\begin{tabular}{ll}
\toprule
\textbf{Stage} & \textbf{Operation} \\
\midrule
0. Pre-projection & Pixel shuffle, factor $0.5$: $1024$ tokens of dim $D_{V} \rightarrow 256$ tokens of dim $4 D_{V}$ \\
1. Norm & LayerNorm$(4 D_{V})$, $\varepsilon = 10^{-5}$ \\
2. Linear-1 & Linear$(4 D_{V} \rightarrow D_{\mathrm{LM}})$, bias \cmark \\
3. Activation & GELU \\
4. Linear-2 & Linear$(D_{\mathrm{LM}} \rightarrow D_{\mathrm{LM}})$, bias \cmark \\
\midrule
\rowcolor{textcolor}\multicolumn{2}{l}{\textit{Per-scale projector parameter count}} \\
Small (1B; $D_{\mathrm{LM}}{=}896$)   & $\sim$4.5M \\
Medium (2B; $D_{\mathrm{LM}}{=}1536$) & $\sim$8.6M \\
Large (4B; $D_{\mathrm{LM}}{=}2048$)  & $\sim$12.6M \\
X-Large (8B; $D_{\mathrm{LM}}{=}3584$) & $\sim$27.5M \\
\bottomrule
\end{tabular}
\end{table}
Following InternVL-2.5/3, we keep depth and activation fixed across scales; only the projector's hidden width tracks the LM. 

\subsection{Language Model Backbones}\label{app:arch-lm}

For the four points on our scaling ladder we use four different sizes from the \textbf{Qwen2.5} family~\citep{qwen25}: $0.5$B, $1.5$B, $3$B, and $7$B parameters. All four share the Qwen2 transformer architecture~\citep{yang2024qwen2technicalreport}---SwiGLU FFN~\citep{shazeer2020glu}, RMSNorm~\citep{zhang2019root}, RoPE position embeddings~\citep{su2024roformer}, grouped-query attention (GQA)~\citep{ainslie2023gqa}, no QK-normalisation. 
They differ only in their depth/width/head budget and in two minor configuration knobs (max position length and embedding tying), summarised in \cref{tab:arch-lm}. 
Unless specified, we always initialise from the \textit{Base} (non-Instruct) checkpoints by default. 

\begin{table}[!h]
\centering
\caption{\textbf{Language-model backbone architecture across the four scaling-ladder scales.} All are Qwen2.5 base checkpoints~\citep{qwen25}. ``Head dim'' is hidden size divided by query head count. Vocabularies are the standard Qwen2.5 tokenizer.}
\label{tab:arch-lm}
\small
\setlength{\tabcolsep}{4pt}
\begin{tabular}{lcccc}
\toprule
\textbf{Scale} & \textbf{Small (1B)} & \textbf{Medium (2B)} & \textbf{Large (4B)} & \textbf{X-Large (8B)} \\
\textbf{Qwen2.5 size} & 0.5B & 1.5B & 3B & 7B \\
\midrule
\rowcolor{textcolor}\multicolumn{5}{l}{\textit{Shared structural choices (all scales)}} \\
Architecture family & \multicolumn{4}{c}{Qwen2 transformer~\citep{yang2024qwen2technicalreport}} \\
Normalisation & \multicolumn{4}{c}{Pre-RMSNorm, $\varepsilon = 10^{-6}$} \\
QK-normalisation & \multicolumn{4}{c}{\xmark} \\
FFN style & \multicolumn{4}{c}{SwiGLU (gated MLP, SiLU activation)} \\
Attention style & \multicolumn{4}{c}{Grouped-query attention (GQA), QKV bias \cmark, O-proj bias \xmark} \\
Positional embeddings & \multicolumn{4}{c}{RoPE, base $\theta = 10^{6}$, no scaling} \\
\midrule
\rowcolor{textcolor}\multicolumn{5}{l}{\textit{Per-scale dimensions}} \\
Layers & 24 & 28 & 36 & 28 \\
Hidden size & 896 & 1536 & 2048 & 3584 \\
Query heads & 14 & 12 & 16 & 28 \\
KV heads (GQA) & 2 & 2 & 2 & 4 \\
Head dim & 64 & 128 & 128 & 128 \\
FFN intermediate & 4{,}864 & 8{,}960 & 11{,}008 & 18{,}944 \\
\midrule
\rowcolor{textcolor}\multicolumn{5}{l}{\textit{Per-scale config knobs}} \\
Max position embedding & 32{,}768 & 131{,}072 & 32{,}768 & 131{,}072 \\
Tied input/output embed. & \cmark & \cmark & \cmark & \xmark \\
Vocabulary size & 151{,}936 & 151{,}936 & 151{,}936 & 152{,}064 \\
LM parameters & $\sim$494M & $\sim$1.54B & $\sim$3.09B & $\sim$7.62B \\
\bottomrule
\end{tabular}
\end{table}

A few cross-scale observations are worth flagging because they surface in our scaling experiments:

\begin{itemize}[leftmargin=*,itemsep=2pt,topsep=2pt]
\item \textbf{Head dimension is not constant.} The $0.5$B model uses $64$-dim heads, while $1.5$B/$3$B/$7$B all use $128$-dim heads. Practitioners scaling pretraining recipes should be aware that the small scale therefore has a slightly different attention behaviour than the rest of the ladder, even though the rest of the architecture is uniform.
\item \textbf{KV-head count is heavily compressed.} GQA ratios are $14{:}2$, $12{:}2$, $16{:}2$, and $28{:}4$ from Small to X-Large---all sub-7B models share the same minimal $2$ KV heads.
\item \textbf{Tied embeddings only at sub-7B.} The $7$B model is the only scale where input/output embeddings are not tied. This costs $\sim$50M extra LM parameters at the X-Large scale.
\item \textbf{Layer count is non-monotonic.} The $3$B model is deeper (36 layers) than the $7$B model (28 layers); $7$B grows primarily by widening (hidden size $2048 \rightarrow 3584$) rather than deepening.
\end{itemize}

These idiosyncrasies are inherited from the Qwen2.5 release and we deliberately do \emph{not} smooth them out, since our purpose is to produce a benchmark whose scaling axis can be reproduced from publicly-released checkpoints rather than to study clean architectural scaling.

\subsection{End-to-end Parameter Accounting}\label{app:arch-totals}

\cref{tab:arch-totals} sums the three components per scale, giving the total trainable-parameter count behind the ``1B / 2B / 4B / 8B'' labels used throughout the paper. The vision encoder and projector together contribute roughly $5$--$60$\% of parameters at the small scale and $\sim$5\% at the X-Large scale.

\begin{table}[!h]
\centering
\caption{\textbf{Total parameter count per scale.} Vision encoder (InternViT-300M-448px-V2.5) is fixed; LM backbone is the corresponding Qwen2.5 size.}
\label{tab:arch-totals}
\small
\begin{tabular}{lcccc}
\toprule
& \textbf{Small} & \textbf{Medium} & \textbf{Large} & \textbf{X-Large} \\
\midrule
Vision encoder & 304M & 304M & 304M & 304M \\
Projector & 4.5M & 8.6M & 12.6M & 27.5M \\
LM backbone & 494M & 1.54B & 3.09B & 7.62B \\
\midrule
\textbf{Total trainable} & \textbf{$\sim$0.80B} & \textbf{$\sim$1.85B} & \textbf{$\sim$3.40B} & \textbf{$\sim$7.95B} \\
Paper label & 1B & 2B & 4B & 8B \\
\bottomrule
\end{tabular}
\end{table}

All parameters are trained jointly---there is no frozen-encoder pretraining stage, no LoRA~\citep{hu2022lora} adapter, and no separate connector-warmup phase. We refer the reader to \cref{tab:hparams} for the optimizer, schedule, and packing settings used during this joint training.
